\clearpage
\section{Strain Tensor}
\subsection{small deformation}
The deformation is 
\begin{equation}
    \vb{u}=\ins{x',y',z'}-\ins{x,y,z}=\ins{x\ins{e^{-y/D}-1},y\ins{e^{-x/D}-1},0}
\end{equation}
for the deformation to be small, we want 
\begin{equation}
    \pdv{u_i}{x_j} \ll 1
\end{equation}
and since the coordinates are of order of magnitude of \(L\) we want 
\begin{equation}
    \frac{L}{D}\ll 1
\end{equation}
or 
\begin{equation}
    L \ll D
\end{equation}
\subsection{nonlinear components}
in the small deformation case we found that 
\begin{equation}
    L/D \ll 1
\end{equation}
and the nonlinear term goes like 
\begin{equation}
    \pdv{u_k}{x_i'}\pdv{u_k}{x_j'} \sim \ins{L/D}^2 \ll L/D
\end{equation}
so it's very negligible.
\subsection{Linear Strain Tensor}

\begin{equation}
    \epsilon_{ij} = \frac{1}{2}\ins{\pdv{u_i}{x_j}+\pdv{u_j}{x_i}}
\end{equation}
where 
\begin{equation}
    \pdv{u_x}{x}\simeq -y/D \qquad \pdv{u_x}{y}\simeq -x/D \qquad \pdv{u_y}{x}\simeq-y/D \qquad \pdv{u_y}{y}\simeq-x/D
\end{equation}
so 
\begin{equation}
    \epsilon_{ij} = \mqty(-y/D&-\frac{x+y}{2D}&0\\-\frac{x+y}{2D}&-x/D&0\\0&0&0)
\end{equation}
\subsection{volume change}
we saw in the lecture that the change in volume of an element is 
\begin{equation}
    \frac{dV}{V}=\epsilon_{ii}
\end{equation}
and here 
\begin{equation}
    \epsilon_{ii}=-y/D-x/D=-\frac{x+y}{D}
\end{equation}
so the total change is 
\begin{equation}
    \Delta V = \int_{o}^{L}\int_{o}^{L}\int_{o}^{L}\ins{-\frac{x+y}{D}}\dd{x}\dd{y}\dd{z}
\end{equation}
which from symmetry (in the mathematical sense) we can rewrite as 
\begin{equation}
    \Delta V = \frac{-2}{D}\int_{o}^{L}\int_{o}^{L}\int_{o}^{L}x\dd{x}\dd{y}\dd{z}=-\frac{L^4}{D}
\end{equation}
so the total volume is 
\begin{equation}
    V' = L^3-\frac{L^4}{D}
\end{equation}